% The left fixed point of a uniform MPS: l E = lambda l, where E is one column
% of the transfer matrix.
%
% E is not drawn as a thing of its own, because it is not one -- it is the ket
% tensor A and its conjugate stacked and joined along the physical index, and
% the picture says that more plainly than the letter does. l carries two bond
% indices, one into each layer, which is why it is drawn tall rather than as a
% site: it lives in the doubled space, not on the chain.
%
% A is left-canonical, so the same triangle as in the chain; the lower row is
% its conjugate and is the same tensor, which is why it is the same shape. The
% physical index between them is the contraction that makes E a map from the
% doubled bond space to itself, and it continues both tensors' own flat sides,
% so the column reads as one object.
%
% Solve it and you have what an infinite chain needs in place of an end: the
% boundary that reproduces itself under one more site.
\documentclass[border=4pt]{standalone}
\usepackage{amsmath,amssymb}
\usepackage{tikz-tensors}
\begin{document}
\begin{tikzpicture}[x=1.3cm]
  % ---- l E --------------------------------------------------------------
  \node[tn node, minimum width=7mm, minimum height=28mm] (l) at (0,0) {$l$};
  \node[canl] (A)  at (1.25,0.85)  {$A$};
  \node[canl] (Ab) at (1.25,-0.85) {$\bar A$};

  \tnbond[disc]{(l.east |- A) -- (A)}
  \tnbond[disc]{(l.east |- Ab) -- (Ab)}
  \tnbond[disc]{(A.corner 3) -- (Ab.corner 2)}
  \tnbond[disc]{(A) -- ++(0.5,0)}
  \tnbond[disc]{(Ab) -- ++(0.5,0)}

  % ---- = lambda l -------------------------------------------------------
  \node at (2.45,0) {$=\;\lambda$};
  \node[tn node, minimum width=7mm, minimum height=28mm] (r) at (3.25,0) {$l$};
  \tnbond[disc]{([yshift=8.5mm]r.east) -- ++(0.5,0)}
  \tnbond[disc]{([yshift=-8.5mm]r.east) -- ++(0.5,0)}
\end{tikzpicture}
\end{document}
